В ходе выполнения практико-ориентированного проекта были получены и закреплены навыки тестирования и контроля качества веб-приложений на примере информационного портала \textit{Цифровая Энергетика}. Разработан комплексный план тестовых мероприятий, управление процессом и трекинг задач успешно организованы с применением облачной системы Kaiten.

Проведено ручное тестирование графического интерфейса: подтверждена кроссбраузерная совместимость в средах Firefox, Opera и Vivaldi, а также корректность адаптивной верстки на мобильных и десктопных разрешениях. Технический аудит и анализ сетевых подключений с помощью встроенных инструментов разработчика и внешнего сервиса W3Techs позволили определить стек технологий и убедиться в безопасности передачи данных. Одной из выявленных проблем стало наличие ошибок валидации HTML-кода по стандарту W3C, что зафиксировано в отчете.

Для перехода к автоматизации проверок были определены CSS- и XPath-локаторы ключевых элементов DOM-дерева, после чего разработан и отлажен сценарий дымового тестирования в Selenium IDE. Интеграция методов искусственного интеллекта позволила оптимизировать проектирование тест-кейсов и теоретическую базу, доказав целесообразность применения LLM в рутинных задачах. Практическая значимость выполненной работы заключается в формировании комплексного подхода к анализу современных веб-систем.